% !TeX program = pdflatex
\documentclass[10pt,a5paper]{article}
\newif\ifmostrarSoluciones
\mostrarSolucionestrue
\usepackage{estilo_apuntes}
\geometry{a4paper,margin=20mm}
\encabezado{Tutoría}{Logaritmos y exponenciales}
\begin{document}
\begin{center}
  {\large\bfseries\color{Azul} Tutoría: $\ln$ y $\exp$}\\[3pt]
  23 de noviembre de 2026 · 60 minutos
\end{center}

\textbf{Objetivo.} Resolver ejercicios de dominio, ecuaciones y límites con $\ln x$ y $\exp(x)=e^x$, justificando las restricciones y los cambios de variable.

\section*{0--8 min: activar lo anterior}
\begin{ejercicios}[repaso breve]
Calcule y explique por qué no basta sustituir directamente:
\[
 \text{a) }\lim_{x\to2}\frac{x^2-4}{x-2},
 \qquad
 \text{b) }\lim_{x\to0}\frac{\sin(3x)}{2x}.
\]
En b), los ángulos están en radianes.
\end{ejercicios}

\section*{8--18 min: dominio y propiedades}
\begin{ejercicios}[antes de operar]
\begin{enumerate}
  \item Determine el dominio de $f(x)=\ln\!\left(\dfrac{x-1}{x+2}\right)$.
  \item Simplifique $e^{\ln(x-1)}$ y $\ln(e^{2x-1})$, indicando el dominio de cada expresión original.
  \item ¿Es cierto que $\ln(a+b)=\ln a+\ln b$ para $a,b>0$? Dé un contraejemplo.
\end{enumerate}
\end{ejercicios}

\section*{18--30 min: resolver y comprobar}
\begin{ejercicios}[ecuaciones]
Resuelva en $\mathbb R$. Escriba primero las restricciones.
\begin{enumerate}
  \item $\ln(x-1)+\ln(x+1)=\ln 8$.
  \item $e^{2x}-3e^x+2=0$. Sugerencia: $u=e^x>0$.
\end{enumerate}
\end{ejercicios}

\section*{30--47 min: límites con $\ln$ y $\exp$}
Recordar únicamente las herramientas que se van a usar:
\[
 \lim_{t\to0}\frac{e^t-1}{t}=1,
 \qquad
 \lim_{t\to0}\frac{\ln(1+t)}{t}=1.
\]
$e^x$ es continua en $\mathbb R$ y $\ln x$ en $(0,\infty)$.
\begin{ejercicios}[de menor a mayor dificultad]
\begin{enumerate}
  \item $\displaystyle\lim_{x\to0}\frac{\ln(2+e^x)}{1+e^x}$.
  \item $\displaystyle\lim_{x\to0}\frac{e^{3x}-1}{2x}$.
  \item $\displaystyle\lim_{x\to0}\frac{\ln(1-2x)}{e^{3x}-1}$.
\end{enumerate}
En cada caso: revisar dominio, intentar sustitución y, si aparece $0/0$, ajustar el cociente al límite fundamental. En el tercero, hacer explícitos los factores $-2$ y $3$.
\end{ejercicios}

\section*{47--55 min: conectar con continuidad}
\begin{ejercicios}[función por tramos]
Determine $k$ para que $F$ sea continua en $0$:
\[
 F(x)=\begin{cases}
 \dfrac{\ln(1+2x)}{e^x-1},&x\ne0,\quad x>-\tfrac12,\\[5pt]
 k,&x=0.
 \end{cases}
\]
¿Qué datos se comparan para justificar la continuidad?
\end{ejercicios}

\section*{55--60 min: propuestos}
Calcule $\displaystyle\lim_{x\to0}\frac{\ln(1+4x)}{e^{2x}-1}$ e indique la restricción del logaritmo. Justifique sin derivar.

\textbf{Reserva (fuera de los 60 minutos).} Calcule
\[
 \lim_{x\to0}\frac{\ln(1+3x)}{\sin(2x)},\qquad
 \lim_{x\to+\infty}\frac{\ln x}{\sqrt{x}},\qquad
 \lim_{x\to+\infty}\frac{x^2}{e^x}.
\]

\ifmostrarSoluciones
\newpage
\section*{Soluciones para dirigir la tutoría}
\textbf{Repaso.} a) Para $x\ne2$, $(x^2-4)/(x-2)=x+2$; el límite es $4$. b) Escribimos $\frac32\frac{\sin(3x)}{3x}$; el límite es $3/2$. En ambos, $0/0$ es una indeterminación, no el resultado.

\textbf{Dominio y propiedades.}
\begin{enumerate}
  \item Se exige $(x-1)/(x+2)>0$. Los puntos críticos son $-2$ y $1$; numerador y denominador tienen el mismo signo en $(-\infty,-2)$ y $(1,\infty)$. Ese es el dominio. No se puede separar en $\ln(x-1)-\ln(x+2)$ en todo el dominio: en $x<-2$ ambos argumentos serían negativos.
  \item $e^{\ln(x-1)}=x-1$ para $x>1$. La simplificación conserva esa restricción. Por otra parte, $\ln(e^{2x-1})=2x-1$ para todo $x\in\mathbb R$, pues la exponencial siempre es positiva.
  \item Con $a=b=1$, el lado izquierdo vale $\ln2$ y el derecho $0$. La propiedad correcta es $\ln(ab)=\ln a+\ln b$ para $a,b>0$.
\end{enumerate}

\textbf{Ecuaciones.}
\begin{enumerate}
  \item Las condiciones $x-1>0$ y $x+1>0$ dan $x>1$. Entonces
  \[
    \ln((x-1)(x+1))=\ln8
    \ \Longrightarrow\ x^2-1=8.
  \]
  Los candidatos son $\pm3$; solo $x=3$ cumple el dominio. Al comprobar: $\ln2+\ln4=\ln8$.
  \item Con $u=e^x>0$, $(u-1)(u-2)=0$. Ambos valores son positivos; por tanto, $x=0$ o $x=\ln2$.
\end{enumerate}

\textbf{Límites.}
\begin{enumerate}
  \item Por continuidad, el límite es $\ln3/2$. El argumento $2+e^x$ y el denominador son positivos para todo $x$.
  \item Para $x\ne0$,
  \[
    \frac{e^{3x}-1}{2x}=\frac32\frac{e^{3x}-1}{3x}
    \longrightarrow\frac32.
  \]
  \item El logaritmo exige $x<1/2$, condición que se cumple en un entorno de $0$. Para $x\ne0$ en ese entorno,
  \[
    \frac{\ln(1-2x)}{e^{3x}-1}
    =\frac{\ln(1-2x)}{-2x}
     \frac{3x}{e^{3x}-1}\left(-\frac23\right)
    \longrightarrow-\frac23.
  \]
  Los argumentos $-2x$ y $3x$ tienden a $0$. El segundo factor tiende a $1$ porque es el recíproco de un cociente cuyo límite es $1\ne0$.
\end{enumerate}

\textbf{Continuidad.} Para $x\ne0$ cercano a $0$,
\[
 \frac{\ln(1+2x)}{e^x-1}
 =2\frac{\ln(1+2x)}{2x}\frac{x}{e^x-1}
 \longrightarrow2.
\]
Los límites izquierdo y derecho valen $2$; deben coincidir con $F(0)=k$. Por tanto, $k=2$. La fórmula del cociente no define el valor en $0$; ese valor se asigna mediante el tramo $k$.

\textbf{Salida.} La restricción es $x>-1/4$. Para $x\ne0$ cercano a $0$,
\[
 \frac{\ln(1+4x)}{e^{2x}-1}
 =2\frac{\ln(1+4x)}{4x}\frac{2x}{e^{2x}-1}
 \longrightarrow2.
\]
\textbf{Reserva.} Los resultados son $3/2$, $0$ y $0$. En el primero se usa
$\frac{\ln(1+3x)}{3x}\frac{2x}{\sin(2x)}\frac32$.
Los otros dos se justifican por el crecimiento comparado visto en semana 7: el logaritmo crece más lentamente que cualquier potencia positiva y la exponencial más rápidamente que cualquier potencia.

\textbf{Criterio de cierre.} El estudiante identifica el dominio, reconoce cuándo puede sustituir y muestra los factores necesarios para aplicar los límites fundamentales. Si hay dificultad con esos factores, retomar el segundo límite antes de avanzar al cociente mixto.
\fi
\end{document}
